\section{Experimental Setup}
\label{sec:setup}

\subsection{Tasks}

We build on \ParEval~\cite{pareval}, the standard benchmark for LLM-generated
parallel code, which comprises 420 problems spanning 12 problem
families---sort, scan, dense and sparse linear algebra, search, reduce,
histogram, stencil, graph, geometry, Fourier transform, and transform---across
seven execution models. We use the OpenMP subset (\NTasks\ tasks), reusing
\ParEval's function-completion prompts, its randomized-input drivers, and most
importantly its serial reference implementations as the trusted answer.
Because the reference is a plain sequential program it is immune to the
parallel hazards under study, which makes it an honest oracle; because the
drivers supply many random inputs, a program cannot pass on one lucky example.

Each prompt states what to compute and gives the function signature, but never
says how to parallelize it. The model chooses the synchronization strategy, and
therefore the scaling behaviour, on its own.

\subsection{Models: who wrote the code}

\NModels\ closed and \NOpenModels\ open-weight models generated every program we
study, and we chose them so that the degeneracy could not be waved away as an
artifact of a weak generator. Four are reached through a single uniform API and
belong to one frontier family, which lets us watch the reward behave across a
capability jump rather than at a single point: \texttt{gpt-5.4}, a current
frontier general model, and \texttt{gpt-5.2}, the generation before it;
\texttt{gpt-5.3-codex}, a code-specialized sibling of exactly the kind an HPC
team reaches for; and \texttt{gpt-5.4-pro}, a higher-reasoning variant that
spends more test-time compute per program. If the reward's blind spot were
something a stronger or more code-aware model simply grows out of, this
panel---two generations, a code specialist, and a reasoning model---is where we
would expect to see it close.

It does not. Every model in the panel, the strongest included, emits programs
whose measured 8-thread speedup ranges from barely above serial to near-linear
and earns full reward for all of them: \texttt{gpt-5.4-pro}, our
higher-reasoning variant, produces accepted programs spanning $1.9\times$ to
$7.4\times$, each scored identically correct. The reward's inability to tell a
$2\times$ program from a $7\times$ one is a property of what it measures, not of
who is answering. To make that point robust to the specific provider, we add an
open-weights panel of
\NOpenModels\ models (\texttt{gpt-oss}, \texttt{minimax-m2},
\texttt{qwen3-small}) served on a research cluster, spanning a wide capability
gradient. Holding the benchmark, harness, and prompts fixed across all seven,
any difference we report is attributable to the reward, not to model style.

\subsection{The frozen sample pool}

All reward comparisons in Section~\ref{sec:results} are computed offline over
a single frozen pool: \NSamples\ generated programs from \NRuns\ runs, each
annotated with its compile status and per-configuration correctness verdict,
and---for the \NAccepted\ programs the verifier accepted---wall time at every
thread count and repetition. Freezing the pool has three
benefits. Every reward formulation is scored on identical samples, so
differences are attributable to the reward and not to decoding. The comparison
costs no inference budget and can be rerun by anyone. And the pool is
releasable in a way that a live API-backed experiment is not.

\subsection{Hardware, compilers, and measurement protocol}
\label{sec:setup-protocol}

Compilation and execution run on a single dedicated multi-core node with no
other tenants, using GCC with \texttt{-O3 -fopenmp}. We deliberately hold the
machine fixed so that every reward comparison is made under identical timing
conditions; because efficiency is architecture-relative, we treat cross-machine
transfer of the efficiency reward as an explicit limitation
(Section~\ref{sec:discussion}) rather than a result.

Timing follows a protocol designed to make the efficiency term reproducible
enough to be a reward rather than an anecdote:

\begin{itemize}
  \item Thread counts $n \in \{1,2,4,\dots,\NThreadsMax\}$, with
  \texttt{OMP\_PROC\_BIND=close} and \texttt{OMP\_PLACES=cores}. Pinning is
  reported because unpinned runs are not comparable across repeats.
  \item \NRepeats\ repetitions per configuration, each in a \emph{fresh
  process} with a freshly randomized input, which closes the cross-run caching
  channel (H7).
  \item The first repetition is discarded as warm-up; we report the minimum
  wall time over the remainder.
  \item The timed region encloses the entire driver call, not a region the
  generated code controls, which closes the timing-region channel (H6).
  \item Input sizes chosen large enough to amortize thread-creation overhead,
  because small inputs systematically hide poor scaling.
\end{itemize}

\subsection{Correctness verification}

Correctness is adjudicated by \ParEval's driver, extended with a differential
sweep: the generated function is compiled against the serial reference, fed
identical random inputs, and checked at every thread count and repetition
rather than once. A program counts as correct only if it agrees everywhere.
Disagreement that varies with thread count, or flickers between repetitions at
a fixed thread count, is reported as the signature of a race or an
order-dependent bug.

As an independent instrument we cross-check every accepted program with
Archer~\cite{archer}, the LLVM happens-before OpenMP race detector built on
ThreadSanitizer~\cite{tsan} with OMPT annotations, which reasons about actual
synchronization rather than observed output. Where the two instruments agree,
the claim rests on more than output differencing.

\subsection{Reward formulations compared}

Table~\ref{tab:rewards} lists the formulations scored over the frozen pool. All
are implemented in the released harness; where a published reward has
hyperparameters we use the values from the original paper and say so.

\begin{table}[t]
\caption{Reward formulations scored over the frozen sample pool.}
\label{tab:rewards}
\centering
\footnotesize
\begin{tabular}{@{}llp{2.7cm}@{}}
\toprule
ID & Form & Source \\
\midrule
R1 & $c$ & correctness-only (the default) \\
R2 & $c(0.5 + 0.5\,\hat{E})$ & ACECode~\cite{acecode} (additive) \\
R3 & $c(0.3 + \hat{E})/1.3$ & Kevin~\cite{kevin} (corr.+speedup) \\
R5 & $c\cdot\min(S_p(n)/n,1)$ & \textbf{this work} (reference-anchored gate) \\
R5b & $c\cdot\min(\Sigma_p(n)/n,1)$ & ablation (self-anchored) \\
R5c & $c\cdot\min(S_p(n)/n,1)\cdot\min(\Sigma_p(n)/2,1)$ & ablation (decomposed) \\
\bottomrule
\end{tabular}
\\[2pt]
{\footnotesize $\hat{E}=\min(S_p(n)/n,1)$; $S_p$ uses the trusted serial
reference, $\Sigma_p$ the program's own one-thread time.}
\end{table}
